\section{从 Chat 到 Agent：范式迁移}

访谈开头把 2026 年定义为大模型战争第二幕。第一幕是 Chat：用户和模型通过问答交互，核心竞争来自预训练规模、指令对齐和短上下文体验。第二幕是 Agent：模型进入环境、调用工具、执行长程任务，核心竞争转向后训练、RL infra、环境反馈、成本和可靠性。

\begin{figure}[H]
\centering
\includegraphics[width=0.94\textwidth]{chat-to-agent-paradigm.png}
\caption{从 Chat/Pre-train 主导到 Agent/Post-train 主导的范式迁移。}
\end{figure}

\begin{knowledgebox}{读图：这张图应该怎么看}
左侧 Chat 阶段强调问答、检索和生成，资源重心在预训练数据与基座模型；右侧 Agent 阶段强调工具、环境和长程任务，资源重心转到 rollout、RL infra、评估和框架。图中箭头不是说预训练不重要，而是说竞争焦点发生转移。
\end{knowledgebox}

\subsection{OpenClaw 的触发作用}

罗福莉最初把 OpenClaw 看作 Claude Code 加一个 UI，本能排斥它的运营包装。但真正使用后，她把它定义为划时代 Agent 框架，因为它让模型短板被框架弥补，也让普通团队成员能通过框架参与“提升智能水平”。OpenClaw 的意义不是某个 UI，而是让 Agent 框架成为研究范式。

\begin{importantbox}{OpenClaw 的本质}
OpenClaw 让模型、工具、任务环境和人的想象力被组织到一个可迭代框架里。它不是单纯“产品形态创新”，而是把后训练时代的研究问题显性化：如何让 Agent 框架和模型互相放大。
\end{importantbox}

\subsection{为什么 2026 是生产力变革之年}

罗福莉认为，中国开发者对 OpenClaw 类工具反应更强，一方面因为效率提升需求更迫切，另一方面因为国内有大量便宜好用模型。若一个复杂 Agent 任务花 10 元 API 成本，却替代 1000 元人力价值，采用动力会非常强。Agent 生产力革命的前置条件，是足够好的框架、足够便宜的模型和足够明确的任务回报。

\begin{warningbox}{不要把火爆等同于成熟}
一个框架火起来，可能说明它抓住了生产力需求，但不代表它已经工业级成熟。真正成熟需要端到端完成率、成本效率、速度、可靠性和安全边界共同达标。
\end{warningbox}

\subsection{本章小结}

Chat 阶段展示模型智能，Agent 阶段要求模型进入任务环境。OpenClaw 的冲击在于，它让许多人第一次感受到：框架、工具和群体使用方式可以快速改变模型能力的外显上限。

